\subsection{Prohibition of Neutrinoless Double Beta Decay (\texorpdfstring{$0\nu\beta\beta$}{0nubetabeta})}

Standard physics often invokes \textbf{Majorana masses} and the \textbf{See-saw Mechanism} to explain the extreme lightness of neutrinos. This requires the existence of yet-undiscovered, ultra-heavy right-handed neutrinos.

The $E_8$-Persistence Theory achieves mass suppression via the \textbf{Inverse Resonance Scaling} ($1/\Delta^n$) derived in Paper II, rendering the See-saw mechanism a mathematical artifact. More importantly, the \textbf{Chiral Diode} architecture (Section IV) imposes a strict topological constraint on the neutrino nature.

\begin{itemize}
    \item \textbf{The Constraint:} The Chiral Diode enforces a unidirectional information flow to preserve causality. A Majorana fermion (which is its own antiparticle) would represent a standing wave in the time dimension, violating the chiral truncation required for the Lorentzian projection. Therefore, neutrinos must be \textbf{Dirac fermions}.
    
    \item \textbf{The Decisive Test:} This represents a ``Hard Falsification" point. The observation of neutrinoless double-beta decay by experiments like LEGEND or nEXO would imply that the Chiral Diode is permeable, falsifying the fundamental geometric logic of this framework.
\end{itemize}